\section{Conclusion}
\label{sec:conclusion}

A constraint inferred from data describes what the data happens to satisfy, whereas a constraint contained in governing text describes what the data is required to satisfy. This distinction has received limited attention in database research, although many published datasets come with documents that define their quality criteria. 

We showed that the gap between data and normative text can be addressed automatically: lightweight classification detects most checkable requirements, and a structured multi-step pipeline grounded in both text and schema translates them into executable constraints while avoiding hallucinated predicates and literals. 

This enables a new form of validation: \system\ can reveal compliance failures even when they do not manifest as discoverable patterns in the data. In such settings, the normative document becomes part of the data specification, and keeping that specification aligned with evolving text emerges as a concrete challenge for data management. Because each constraint comes from the text, it remains valid while the regulation is unchanged and can be re-executed, also incrementally~\cite{youri2024}, as the data evolves; only a change in the text requires a new mapping. Legal documents often have hundreds of sentences with nested clauses that must first be segmented and resolved, against at most 158 sentences in our datasets; since cost grows linearly with the number of sentences, we expect the pipeline to scale, but this remains to be shown.


